% ============================================================
%  CAHIER DES DIAGRAMMES UML — APPLICATION CAYA
%  Club des Amis de Yaoundé
%  Version 1.0 — Mars 2026
% ============================================================
\documentclass[12pt, a4paper]{report}

\usepackage[utf8]{inputenc}
\usepackage[T1]{fontenc}
\usepackage[french]{babel}
\usepackage[top=2.5cm, bottom=2.5cm, left=2.8cm, right=2.8cm]{geometry}
\usepackage{setspace}
\onehalfspacing
\usepackage{lmodern}
\usepackage{microtype}

% Couleurs
\usepackage[dvipsnames, table]{xcolor}
\definecolor{cayaBlue}{RGB}{0, 71, 130}
\definecolor{cayaGold}{RGB}{201, 155, 0}
\definecolor{cayaLight}{RGB}{230, 241, 255}
\definecolor{warningRed}{RGB}{200, 40, 40}
\definecolor{successGreen}{RGB}{30, 130, 60}
\definecolor{codeGray}{RGB}{248, 248, 248}
\definecolor{plantKeyword}{RGB}{0, 0, 180}
\definecolor{plantComment}{RGB}{60, 128, 60}
\definecolor{plantString}{RGB}{160, 32, 32}
\definecolor{plantType}{RGB}{128, 0, 128}

% En-têtes
\usepackage{fancyhdr}
\pagestyle{fancy}
\fancyhf{}
\fancyhead[L]{\small\textcolor{cayaBlue}{\textbf{CAYA} — Cahier des Diagrammes UML}}
\fancyhead[R]{\small\textcolor{gray}{v1.0 · Mars 2026}}
\fancyfoot[C]{\small\thepage}
\renewcommand{\headrulewidth}{0.4pt}

% Tableaux
\usepackage{booktabs}
\usepackage{tabularx}
\usepackage{array}
\usepackage{colortbl}
\newcolumntype{L}[1]{>{\raggedright\arraybackslash}p{#1}}
\newcolumntype{C}[1]{>{\centering\arraybackslash}p{#1}}

% Graphiques
\usepackage{graphicx}
\usepackage{tikz}
\usetikzlibrary{shapes.geometric, arrows.meta, positioning, fit, backgrounds, calc}

% Listes
\usepackage{enumitem}
\setlist[itemize]{leftmargin=1.5em, itemsep=2pt}

% Boîtes
\usepackage{tcolorbox}
\tcbuselibrary{breakable, skins, listings}

\newtcolorbox{diagramBox}[2]{
  colback=codeGray, colframe=cayaBlue,
  fonttitle=\bfseries\small\color{white},
  title=\textbf{#1} \hfill {\normalfont\small #2},
  coltitle=white, attach boxed title to top left,
  boxed title style={colback=cayaBlue, arc=2pt},
  breakable, boxrule=0.8pt, arc=3pt,
  before upper={\setlength{\parindent}{0pt}}
}

\newtcolorbox{noteBox}[1]{
  colback=cayaLight, colframe=cayaBlue!60,
  fonttitle=\bfseries\small, title=#1,
  breakable, boxrule=0.6pt, arc=3pt
}

% Code PlantUML
\usepackage{listings}
\lstdefinelanguage{PlantUML}{
  keywords={@startuml,@enduml,class,interface,enum,abstract,
            participant,actor,boundary,control,entity,database,
            note,end,left,right,of,on,link,as,
            skinparam,title,package,namespace,rectangle,node,
            usecase,component,artifact,folder,frame,
            state,choice,fork,join,[*],
            activate,deactivate,destroy,
            autonumber,ref,over,loop,alt,else,opt,
            break,critical,group,par,strict},
  keywordstyle=\color{plantKeyword}\bfseries,
  comment=[l]{'},
  commentstyle=\color{plantComment}\itshape,
  stringstyle=\color{plantString},
  morestring=[b]",
  sensitive=false,
  moredelim=[s][\color{plantType}]{<}{>},
}
\lstset{
  language=PlantUML,
  basicstyle=\ttfamily\scriptsize,
  backgroundcolor=\color{codeGray},
  frame=none,
  breaklines=true,
  keepspaces=true,
  showstringspaces=false,
  numbers=left,
  numberstyle=\tiny\color{gray!60},
  numbersep=5pt,
  xleftmargin=10pt,
  tabsize=2
}

% Liens
\usepackage[hidelinks, colorlinks=true, linkcolor=cayaBlue]{hyperref}
\usepackage{bookmark}
\usepackage{pifont}
\newcommand{\cmark}{\textcolor{successGreen}{\ding{51}}}

% ============================================================
\title{%
  \vspace{-1cm}
  {\Huge\bfseries\color{cayaBlue} Cahier des Diagrammes UML}\\[0.4cm]
  {\large\bfseries Application de Gestion\\[0.1cm]
  \textbf{\color{cayaGold}CLUB DES AMIS DE YAOUNDÉ (CAYA)}}\\[0.3cm]
  {\normalsize 18 Diagrammes · PlantUML · Classes · Séquences · États · Déploiement}\\[0.5cm]
  \rule{0.6\textwidth}{1pt}\\[0.4cm]
  {\small Version 1.0 \quad|\quad Mars 2026 \quad|\quad Confidentiel}
}
\author{}
\date{}

% ============================================================
\begin{document}
\maketitle
\thispagestyle{empty}

\vfill
\begin{center}
\begin{tabular}{ll}
  \toprule
  \textbf{Outil de rendu} & PlantUML (\texttt{https://www.plantuml.com/plantuml}) \\
  \textbf{Modules couverts} & 12 modules métier \\
  \textbf{Diagrammes produits} & 18 (Classes × 11, États × 4, Séquences × 3) \\
  \textbf{Prérequis} & Cahier des Charges + Cahier de Conception v1.0 \\
  \bottomrule
\end{tabular}
\end{center}

\newpage
\tableofcontents
\newpage

% ============================================================
\chapter*{Guide de Lecture}
\addcontentsline{toc}{chapter}{Guide de Lecture}
% ============================================================

\begin{noteBox}{Comment utiliser ce document}
Chaque bloc de code est du \textbf{PlantUML} valide. Pour le rendre :
\begin{enumerate}
  \item Copier le contenu entre \texttt{@startuml} et \texttt{@enduml}
  \item Le coller sur \texttt{https://www.plantuml.com/plantuml/uml/}
  \item Ou utiliser l'extension \textit{PlantUML} dans VSCode avec \texttt{Alt+D}
\end{enumerate}
Chaque diagramme est autonome et peut être généré indépendamment.
\end{noteBox}

\begin{center}
\rowcolors{2}{cayaLight}{white}
\begin{tabular}{C{0.8cm}L{3.5cm}L{2.5cm}L{6cm}}
  \toprule
  \rowcolor{cayaBlue}
  \textcolor{white}{\textbf{\#}} &
  \textcolor{white}{\textbf{Diagramme}} &
  \textcolor{white}{\textbf{Type}} &
  \textcolor{white}{\textbf{Description}} \\
  \midrule
  D01 & Auth \& RBAC             & Classes    & Utilisateurs, rôles, tokens \\
  D02 & Configuration            & Classes    & TontineConfig, snapshot exercice \\
  D03 & Membres                  & Classes    & Profils, contacts, parrainage \\
  D04 & Exercice Fiscal          & Classes    & FiscalYear, Membership, Parts \\
  D05 & Sessions \& Transactions & Classes    & Sessions, entrées, séquences \\
  D06 & Épargne \& Intérêts      & Classes    & Ledgers, redistribution, pool \\
  D07 & Prêts                    & Classes    & LoanAccount, accrual, rembt \\
  D08 & Caisse de Secours        & Classes    & RescueFund, événements \\
  D09 & Bénéficiaires            & Classes    & Planning, slots, désignation \\
  D10 & Cassation \& Audit       & Classes    & Clôture, rollover, AuditLog \\
  D11 & Vue Globale              & Classes    & Relations inter-modules \\
  D12 & Cas d'Utilisation        & Use Case   & Acteurs et fonctions \\
  D13 & Composants               & Composants & Architecture technique \\
  D14 & Déploiement              & Déploiement& Infrastructure serveurs \\
  D15 & État — FiscalYear        & États      & Cycle de vie exercice \\
  D16 & État — MonthlySession    & États      & Cycle de vie session \\
  D17 & État — LoanAccount       & États      & Cycle de vie prêt \\
  D18 & État — MemberStatus      & États      & Statuts membres \\
  D19 & Séquence — Session       & Séquence   & Enregistrement réunion mensuelle \\
  D20 & Séquence — Intérêts      & Séquence   & Distribution mensuelle \\
  D21 & Séquence — Cassation     & Séquence   & Processus de clôture exercice \\
  \bottomrule
\end{tabular}
\end{center}

% ============================================================
\chapter{Diagrammes de Classes}
% ============================================================

% -----------------------------------------------------------
\section{D01 — Module Auth \& RBAC}
% -----------------------------------------------------------

\begin{diagramBox}{D01 — Auth \& RBAC}{Diagramme de Classes}
\begin{lstlisting}
@startuml D01_Auth_RBAC
!theme plain
skinparam classAttributeIconSize 0
skinparam classFontSize 11
skinparam classBackgroundColor #EEF4FF
skinparam classBorderColor #004782
skinparam arrowColor #004782
skinparam packageBackgroundColor #F8F8F8
title <b>Module 01 — Authentification & Contrôle d'Accès (RBAC)</b>

enum BureauRole {
  SUPER_ADMIN
  PRESIDENT
  VICE_PRESIDENT
  TRESORIER
  TRESORIER_ADJOINT
  SECRETAIRE_GENERAL
  SECRETAIRE_ADJOINT
  COMMISSAIRE_AUX_COMPTES
  MEMBRE
}

class User {
  + id            : VARCHAR(10) <<PK>>
  + username      : VARCHAR(50) <<UNIQUE>>
  + phone         : VARCHAR(20) <<UNIQUE>>
  + passwordHash  : VARCHAR(255)
  + role          : BureauRole
  + isActive      : Boolean = true
  + createdAt     : Timestamp
  + updatedAt     : Timestamp
  + lastLoginAt   : Timestamp
  --
  + activate()    : void
  + deactivate()  : void
  + changeRole(r) : void
}

class RefreshToken {
  + id         : VARCHAR(10) <<PK>>
  + userId     : VARCHAR(10) <<FK>>
  + tokenHash  : VARCHAR(255) <<UNIQUE>>
  + expiresAt  : Timestamp
  + revokedAt  : Timestamp
  + createdAt  : Timestamp
  + userAgent  : VARCHAR(500)
  --
  + isValid()  : Boolean
  + revoke()   : void
}

class AuditLog {
  + id         : VARCHAR(10) <<PK>>
  + actorId    : VARCHAR(10) <<FK>>
  + action     : VARCHAR(100)
  + entityType : VARCHAR(50)
  + entityId   : VARCHAR(10)
  + beforeData : JSON
  + afterData  : JSON
  + reason     : TEXT
  + ipAddress  : VARCHAR(45)
  + createdAt  : Timestamp
  --
  <<PARTITION BY YEAR(createdAt)>>
  <<APPEND-ONLY — pas de UPDATE/DELETE>>
}

note right of AuditLog
  Lecture restreinte :
  SUPER_ADMIN
  PRESIDENT
  SECRETAIRE_GENERAL
end note

note bottom of User
  INV-AUTH-01 : 1 seul PRESIDENT actif
  INV-AUTH-02 : 1 seul VICE_PRESIDENT actif
  INV-AUTH-03 : phone synchro MemberProfile.phone1
  INV-AUTH-04 : changeRole() → AuditLog obligatoire
end note

User "1" *-- "0..*" RefreshToken : possède >
User "1" *-- "0..*" AuditLog     : effectue >
User --> BureauRole               : a le rôle >

@enduml
\end{lstlisting}
\end{diagramBox}

% -----------------------------------------------------------
\section{D02 — Module Configuration}
% -----------------------------------------------------------

\begin{diagramBox}{D02 — Configuration}{Diagramme de Classes}
\begin{lstlisting}
@startuml D02_Configuration
!theme plain
skinparam classAttributeIconSize 0
skinparam classFontSize 10
skinparam classBackgroundColor #EEF4FF
skinparam classBorderColor #004782
skinparam arrowColor #004782
title <b>Module 02 — Configuration CAYA</b>

class TontineConfig <<Singleton>> {
  + id                       : VARCHAR = 'caya' <<PK>>
  .. Identité ..
  + name                     : VARCHAR(255)
  + acronym                  : VARCHAR(20)
  + foundedYear              : SMALLINT
  + motto                    : VARCHAR(500)
  + headquartersCity         : VARCHAR(255)
  + registrationNumber       : VARCHAR(100)
  .. Parts ..
  + shareUnitAmount          : DECIMAL(15,2) = 100000
  + halfShareAmount          : DECIMAL(15,2) = 50000
  + potMonthlyAmount         : DECIMAL(15,2) = 3000
  + maxSharesPerMember       : TINYINT = 10
  .. Épargne ..
  + mandatoryInitialSavings  : DECIMAL(15,2) = 100000
  .. Prêts ..
  + loanMonthlyRate          : DECIMAL(5,4) = 0.0400
  + minLoanAmount            : DECIMAL(15,2) = 10000
  + maxLoanAmount            : DECIMAL(15,2) = 1000000
  + maxLoanMultiplier        : TINYINT = 5
  + minSavingsToLoan         : DECIMAL(15,2) = 100000
  + maxConcurrentLoans       : TINYINT = 2
  .. Secours ..
  + rescueFundTarget         : DECIMAL(15,2) = 50000
  + rescueFundMinBalance     : DECIMAL(15,2) = 25000
  .. Inscription ..
  + registrationFeeNew       : DECIMAL(15,2)
  + registrationFeeReturning : DECIMAL(15,2)
  .. Méta ..
  + updatedAt                : Timestamp
  + updatedById              : VARCHAR(10) <<FK>>
}

class FiscalYearConfig <<Snapshot Immuable>> {
  + id                      : VARCHAR(10) <<PK>>
  + fiscalYearId            : VARCHAR(10) <<FK, UNIQUE>>
  + snapshotAt              : Timestamp
  + snapshotById            : VARCHAR(10) <<FK>>
  .. Règles gelées (copie TontineConfig) ..
  + shareUnitAmount         : DECIMAL(15,2)
  + loanMonthlyRate         : DECIMAL(5,4)
  + maxLoanMultiplier       : TINYINT
  + minSavingsToLoan        : DECIMAL(15,2)
  + maxConcurrentLoans      : TINYINT
  + rescueFundTarget        : DECIMAL(15,2)
  + rescueFundMinBalance    : DECIMAL(15,2)
  + registrationFeeNew      : DECIMAL(15,2)
  + registrationFeeReturning: DECIMAL(15,2)
  + interestPoolMethod      : InterestPoolMethod
  .. Exception SUPER_ADMIN ..
  + forcedModifiedAt        : Timestamp
  + forcedModifiedById      : VARCHAR(10)
  + forcedModifiedReason    : TEXT
}

enum InterestPoolMethod {
  THEORETICAL
  ACTUAL
}

class RescueEventAmount {
  + eventType  : RescueEventType <<PK>>
  + amount     : DECIMAL(15,2)
  + label      : VARCHAR(255)
  + updatedById: VARCHAR(10) <<FK>>
  + updatedAt  : Timestamp
}

enum RescueEventType {
  MEMBER_DEATH   = 300000
  RELATIVE_DEATH = 100000
  MARRIAGE       = 200000
  BIRTH          = 25000
  ILLNESS        = 50000
  PROMOTION      = 50000
}

note right of FiscalYearConfig
  CONF-01 : Créé au moment de activate()
  CONF-02 : Immuable après création
  Exception : SUPER_ADMIN peut forcer
              + forcedModifiedReason requis
              + AuditLog obligatoire
  CONF-03 : interestPoolMethod
            switchable avant ACTIVE seulement
end note

TontineConfig    ..> FiscalYearConfig  : snapshot() au ACTIVE
FiscalYearConfig --> InterestPoolMethod
RescueEventAmount --> RescueEventType

@enduml
\end{lstlisting}
\end{diagramBox}

% -----------------------------------------------------------
\section{D03 — Module Membres}
% -----------------------------------------------------------

\begin{diagramBox}{D03 — Membres}{Diagramme de Classes}
\begin{lstlisting}
@startuml D03_Membres
!theme plain
skinparam classAttributeIconSize 0
skinparam classFontSize 11
skinparam classBackgroundColor #EEF4FF
skinparam classBorderColor #004782
skinparam arrowColor #004782
title <b>Module 03 — Gestion des Membres</b>

class User {
  + id       : VARCHAR(10) <<PK>>
  + username : VARCHAR(50)
  + phone    : VARCHAR(20)
  + role     : BureauRole
  + isActive : Boolean
}

class MemberProfile {
  + id                 : VARCHAR(10) <<PK>>
  + memberCode         : VARCHAR(10) <<UNIQUE>>
  + userId             : VARCHAR(10) <<FK, UNIQUE>>
  + firstName          : VARCHAR(100)
  + lastName           : VARCHAR(100)
  + phone1             : VARCHAR(20)
  + phone2             : VARCHAR(20)
  + attachmentPhoto    : MEDIUMBLOB
  + neighborhood       : VARCHAR(255)
  + locationDetail     : TEXT
  + mobileMoneyType    : VARCHAR(50)
  + mobileMoneyNumber  : VARCHAR(20)
  + sponsorId          : VARCHAR(10) <<FK>>
  + createdAt          : Timestamp
  + updatedAt          : Timestamp
  --
  + fullName()         : String
  + getActiveMembership() : Membership
}

class EmergencyContact {
  + id        : VARCHAR(10) <<PK>>
  + profileId : VARCHAR(10) <<FK>>
  + fullName  : VARCHAR(255)
  + phone     : VARCHAR(20)
  + relation  : VARCHAR(100)
}

class MemberCodeSequence <<Singleton>> {
  - lastCode : INT = 0
  --
  + next()   : String
  ' format retourné : CAYA-00042
  ' SELECT ... FOR UPDATE (atomique)
}

note right of MemberProfile
  memberCode : auto-généré
  Format : CAYA-NNNNN
  via MemberCodeSequence.next()

  attachmentPhoto : JPEG compressé
  Stocké en MEDIUMBLOB (max 100 KB)
  — pas d'URL externe —

  sponsorId : parrain obligatoire
  à l'inscription (parrainage)
end note

User          "1"    --  "1"    MemberProfile    : correspond à >
MemberProfile "1"    *-- "1..*" EmergencyContact : possède >
MemberProfile "0..1" --> "0..*" MemberProfile    : parraine >
MemberProfile ..> MemberCodeSequence             : génère code via >

@enduml
\end{lstlisting}
\end{diagramBox}

% -----------------------------------------------------------
\section{D04 — Exercice Fiscal et Adhésions}
% -----------------------------------------------------------

\begin{diagramBox}{D04 — Exercice Fiscal \& Adhésions}{Diagramme de Classes}
\begin{lstlisting}
@startuml D04_Exercice
!theme plain
skinparam classAttributeIconSize 0
skinparam classFontSize 10
skinparam classBackgroundColor #EEF4FF
skinparam classBorderColor #004782
skinparam arrowColor #004782
title <b>Module 04 — Exercice Fiscal & Adhésions</b>

enum FiscalYearStatus {
  PENDING
  ACTIVE
  CASSATION
  CLOSED
  ARCHIVED
}

enum MemberStatus {
  ACTIVE
  INACTIVE
  SUSPENDED
  EXCLUDED
}

enum EnrollmentType {
  NEW
  RETURNING
  MID_YEAR
}

class FiscalYear <<AggregateRoot>> {
  + id             : VARCHAR(10) <<PK>>
  + label          : VARCHAR(20) <<UNIQUE>>
  + startDate      : DATE
  + endDate        : DATE
  + cassationDate  : DATE
  + loanDueDate    : DATE
  + status         : FiscalYearStatus = PENDING
  + openedAt       : Timestamp
  + openedById     : VARCHAR(10) <<FK>>
  + closedAt       : Timestamp
  + closedById     : VARCHAR(10) <<FK>>
  + notes          : TEXT
  --
  + activate()     : void
  + openCassation(): void
  + close()        : void
  + archive()      : void
}

class FiscalYearConfig {
  + id                  : VARCHAR(10) <<PK>>
  + fiscalYearId        : VARCHAR(10) <<FK, UNIQUE>>
  + interestPoolMethod  : InterestPoolMethod
  + shareUnitAmount     : DECIMAL(15,2)
  + loanMonthlyRate     : DECIMAL(5,4)
  ... (snapshot complet TontineConfig)
}

class Membership {
  + id                   : VARCHAR(10) <<PK>>
  + profileId            : VARCHAR(10) <<FK>>
  + fiscalYearId         : VARCHAR(10) <<FK>>
  + status               : MemberStatus = ACTIVE
  + joinedAt             : DATE
  + joinedAtMonth        : TINYINT
  + enrollmentType       : EnrollmentType
  + catchUpAmount        : DECIMAL(15,2) = 0
  + catchUpPaid          : Boolean = false
  + registrationFeePaid  : Boolean = false
  + rescueContribPaid    : Boolean = false
  + initialSavingsPaid   : Boolean = false
  --
  + isEligibleForLoan()  : Boolean
  + computeCatchUp()     : Decimal
  <<UNIQUE(profileId, fiscalYearId)>>
}

class ShareCommitment {
  + id             : VARCHAR(10) <<PK>>
  + membershipId   : VARCHAR(10) <<FK, UNIQUE>>
  + sharesCount    : DECIMAL(4,2)
  + monthlyAmount  : DECIMAL(15,2)
  + isLocked       : Boolean = false
  + lockedAt       : Timestamp
  + lockedById     : VARCHAR(10) <<FK>>
  --
  + lock()          : void
  + getMonthlyDue() : Decimal
}

note right of FiscalYear
  CHECK :
  startDate < loanDueDate
           < cassationDate
           <= endDate

  TRIGGER MySQL :
  Aucun chevauchement de dates
  avec un exercice existant
end note

note bottom of ShareCommitment
  monthlyAmount = sharesCount × shareUnitAmount
  Figé à l'inscription — immuable ensuite

  CHECK : sharesCount >= 0.25
  CHECK : sharesCount <= 10.00
  CHECK : (sharesCount × 100) MOD 25 = 0
  Valeurs valides : 0.25, 0.50, 0.75, 1.00 ... 10.00
end note

note bottom of Membership
  catchUpAmount =
    (joinedAtMonth - 1)
    × (sharesCount × shareUnitAmount
       + rescueFundMonthlyContrib)
  Si enrollmentType = MID_YEAR
end note

FiscalYear  "1"   *-- "0..*" Membership        : contient >
FiscalYear  "1"   *-- "1"    FiscalYearConfig   : configure >
Membership  "1"   *-- "0..1" ShareCommitment    : engage >
FiscalYear  --> FiscalYearStatus
Membership  --> MemberStatus
Membership  --> EnrollmentType

@enduml
\end{lstlisting}
\end{diagramBox}

% -----------------------------------------------------------
\section{D05 — Sessions et Transactions}
% -----------------------------------------------------------

\begin{diagramBox}{D05 — Sessions \& Transactions}{Diagramme de Classes}
\begin{lstlisting}
@startuml D05_Sessions
!theme plain
skinparam classAttributeIconSize 0
skinparam classFontSize 10
skinparam classBackgroundColor #EEF4FF
skinparam classBorderColor #004782
skinparam arrowColor #004782
title <b>Module 05 — Sessions Mensuelles & Transactions</b>

enum SessionStatus {
  DRAFT
  OPEN
  REVIEWING
  CLOSED
}

enum TransactionType {
  INSCRIPTION
  SECOURS
  COTISATION
  POT
  RBT_PRINCIPAL
  RBT_INTEREST
  EPARGNE
  PROJET
  AUTRES
}

class MonthlySession {
  + id                 : VARCHAR(10) <<PK>>
  + fiscalYearId       : VARCHAR(10) <<FK>>
  + sessionNumber      : TINYINT
  + meetingDate        : DATE
  + location           : VARCHAR(500)
  + hostMembershipId   : VARCHAR(10) <<FK>>
  + status             : SessionStatus = DRAFT
  + openedAt           : Timestamp
  + openedById         : VARCHAR(10) <<FK>>
  + closedAt           : Timestamp
  + closedById         : VARCHAR(10) <<FK>>
  .. Totaux snapshot (calculés à la clôture) ..
  + totalInscription   : DECIMAL(15,2) = 0
  + totalSecours       : DECIMAL(15,2) = 0
  + totalCotisation    : DECIMAL(15,2) = 0
  + totalPot           : DECIMAL(15,2) = 0
  + totalRbtPrincipal  : DECIMAL(15,2) = 0
  + totalRbtInterest   : DECIMAL(15,2) = 0
  + totalEpargne       : DECIMAL(15,2) = 0
  + totalProjet        : DECIMAL(15,2) = 0
  + totalAutres        : DECIMAL(15,2) = 0
  --
  + open(byId)         : void
  + requestReview()    : void
  + close(byId)        : void
  + computeTotals()    : void
  <<UNIQUE(fiscalYearId, sessionNumber)>>
}

class SessionEntry {
  + id                : VARCHAR(10) <<PK>>
  + reference         : VARCHAR(30) <<UNIQUE>>
  + sessionId         : VARCHAR(10) <<FK>>
  + membershipId      : VARCHAR(10) <<FK>>
  + type              : TransactionType
  + amount            : DECIMAL(15,2)
  + loanId            : VARCHAR(10) <<FK>>
  + isOutOfSession    : Boolean = false
  + outOfSessionAt    : Timestamp
  + outOfSessionRef   : VARCHAR(100)
  + isImported        : Boolean = false
  + notes             : TEXT
  + recordedById      : VARCHAR(10) <<FK>>
  + recordedAt        : Timestamp
}

class TransactionSequence {
  + fiscalYearId  : VARCHAR(10) <<PK>>
  + sessionNumber : TINYINT     <<PK>>
  + txType        : VARCHAR(20) <<PK>>
  + lastSequence  : INT = 0
  --
  + nextRef()     : String
  ' CAYA-AAAA-MM-TYP-NNNN
  ' SELECT FOR UPDATE (atomique)
}

note right of SessionEntry
  SESS-01 : SessionEntry ne peut être créée
    que si session.status = OPEN
    (sauf isOutOfSession = true)

  SESS-02 : Si isOutOfSession = true :
    type ∈ {RBT_PRINCIPAL, RBT_INTEREST}
    sessionId = NULL

  SESS-03 : Remboursement prêt
    = 2 entries atomiques (même loanId) :
    - 1 RBT_INTEREST
    - 1 RBT_PRINCIPAL

  Trigger MySQL :
  BEFORE INSERT — vérifie session.status
end note

note right of TransactionSequence
  Format référence :
  CAYA-2026-03-COT-0042
       AAAA MM  TYP NNNN

  Types courts :
  INS, SEC, COT, POT,
  RIN, RPR, EPG, PRJ, AUT
end note

MonthlySession "1" *-- "0..*" SessionEntry       : enregistre >
SessionEntry   ..> TransactionSequence            : génère ref via >
SessionEntry   --> TransactionType
MonthlySession --> SessionStatus

@enduml
\end{lstlisting}
\end{diagramBox}

% -----------------------------------------------------------
\section{D06 — Épargne et Redistribution des Intérêts}
% -----------------------------------------------------------

\begin{diagramBox}{D06 — Épargne \& Intérêts}{Diagramme de Classes}
\begin{lstlisting}
@startuml D06_Epargne
!theme plain
skinparam classAttributeIconSize 0
skinparam classFontSize 10
skinparam classBackgroundColor #EEF4FF
skinparam classBorderColor #004782
skinparam arrowColor #004782
title <b>Module 06 — Épargne & Redistribution des Intérêts</b>

enum SavingsEntryType {
  DEPOSIT
  INTEREST_CREDIT
}

enum PoolParticipantType {
  RESCUE_FUND
  BUREAU
}

class SavingsLedger {
  + id                      : VARCHAR(10) <<PK>>
  + membershipId            : VARCHAR(10) <<FK, UNIQUE>>
  + balance                 : DECIMAL(15,2) = 0
  + principalBalance        : DECIMAL(15,2) = 0
  + totalInterestReceived   : DECIMAL(15,2) = 0
  + lastUpdatedAt           : Timestamp
  --
  + deposit(amount, session)  : void
  + creditInterest(amount, snap) : void
  + getNAP()                  : Decimal
}

class SavingsEntry {
  + id             : VARCHAR(10) <<PK>>
  + ledgerId       : VARCHAR(10) <<FK>>
  + sessionId      : VARCHAR(10) <<FK>>
  + sessionEntryId : VARCHAR(10) <<FK>>
  + month          : TINYINT
  + amount         : DECIMAL(15,2)
  + type           : SavingsEntryType
  + balanceAfter   : DECIMAL(15,2)
  + createdAt      : Timestamp
}

class PoolParticipant {
  + id                     : VARCHAR(10) <<PK>>
  + fiscalYearId           : VARCHAR(10) <<FK>>
  + type                   : PoolParticipantType
  + label                  : VARCHAR(100)
  + initialBalance         : DECIMAL(15,2)
  + currentBalance         : DECIMAL(15,2) = 0
  + totalInterestReceived  : DECIMAL(15,2) = 0
  --
  <<UNIQUE(fiscalYearId, type)>>
}

class InterestDistributionSnapshot {
  + id                 : VARCHAR(10) <<PK>>
  + sessionId          : VARCHAR(10) <<FK, UNIQUE>>
  + totalInterestPool  : DECIMAL(15,2)
  + totalSavingsBase   : DECIMAL(15,2)
  + distributedAt      : Timestamp
  + executedById       : VARCHAR(10) <<FK>>
  --
  + monthlyYield()     : Decimal
  ' = totalInterestPool / totalSavingsBase
}

class InterestAllocation {
  + id              : VARCHAR(10) <<PK>>
  + snapshotId      : VARCHAR(10) <<FK>>
  + membershipId    : VARCHAR(10) <<FK>>
  + savingsBalance  : DECIMAL(15,2)
  + allocationAmount: DECIMAL(15,2)
  --
  <<UNIQUE(snapshotId, membershipId)>>
}

note right of InterestDistributionSnapshot
  Formule de redistribution :
  ─────────────────────────────
  totalSavingsBase =
    Σ SavingsLedger.balance
    + Σ PoolParticipant.currentBalance

  Si THEORETICAL :
    pool = Σ MonthlyLoanAccrual.interestAccrued
  Si ACTUAL :
    pool = Σ SessionEntry[RBT_INTEREST].amount

  allocationAmount[i] =
    balance[i] / totalSavingsBase × pool

  Invariant :
  Σ allocationAmount = totalInterestPool ± 1
end note

SavingsLedger "1" *-- "0..*" SavingsEntry               : historise >
InterestDistributionSnapshot "1" *-- "0..*" InterestAllocation : détaille >
SavingsEntry --> SavingsEntryType
PoolParticipant --> PoolParticipantType

@enduml
\end{lstlisting}
\end{diagramBox}

% -----------------------------------------------------------
\section{D07 — Prêts et Remboursements}
% -----------------------------------------------------------

\begin{diagramBox}{D07 — Prêts \& Remboursements}{Diagramme de Classes}
\begin{lstlisting}
@startuml D07_Prets
!theme plain
skinparam classAttributeIconSize 0
skinparam classFontSize 10
skinparam classBackgroundColor #EEF4FF
skinparam classBorderColor #004782
skinparam arrowColor #004782
title <b>Module 07 — Prêts & Remboursements (Encours Glissant)</b>

enum LoanStatus {
  PENDING
  ACTIVE
  PARTIALLY_REPAID
  CLOSED
}

class LoanAccount <<AggregateRoot>> {
  + id                    : VARCHAR(10) <<PK>>
  + membershipId          : VARCHAR(10) <<FK>>
  + fiscalYearId          : VARCHAR(10) <<FK>>
  + principalAmount       : DECIMAL(15,2)
  + monthlyRate           : DECIMAL(5,4) = 0.0400
  + disbursedAt           : DATE
  + disbursedById         : VARCHAR(10) <<FK>>
  + dueBeforeDate         : DATE
  + status                : LoanStatus = PENDING
  + outstandingBalance    : DECIMAL(15,2) = 0
  + totalInterestAccrued  : DECIMAL(15,2) = 0
  + totalRepaid           : DECIMAL(15,2) = 0
  + requestedAt           : Timestamp
  + requestNotes          : TEXT
  --
  + disburse(byId)        : void
  + applyRepayment(amount): void
  + close()               : void
}

class MonthlyLoanAccrual {
  + id                    : VARCHAR(10) <<PK>>
  + loanId                : VARCHAR(10) <<FK>>
  + sessionId             : VARCHAR(10) <<FK>>
  + month                 : TINYINT
  + balanceAtMonthStart   : DECIMAL(15,2)
  + interestAccrued       : DECIMAL(15,2)
  + balanceWithInterest   : DECIMAL(15,2)
  + repaymentReceived     : DECIMAL(15,2) = 0
  + balanceAtMonthEnd     : DECIMAL(15,2)
  --
  <<UNIQUE(loanId, sessionId)>>
  CHECK : balanceAtMonthEnd >= 0
}

class LoanRepayment {
  + id             : VARCHAR(10) <<PK>>
  + loanId         : VARCHAR(10) <<FK>>
  + sessionId      : VARCHAR(10) <<FK>>
  + sessionEntryId : VARCHAR(10) <<FK>>
  + amount         : DECIMAL(15,2)
  + principalPart  : DECIMAL(15,2)
  + interestPart   : DECIMAL(15,2)
  + balanceAfter   : DECIMAL(15,2)
  + recordedAt     : Timestamp
}

class CarryoverLoanRecord {
  + id               : VARCHAR(10) <<PK>>
  + originalLoanId   : VARCHAR(10) <<FK, UNIQUE>>
  + newFiscalYearId  : VARCHAR(10) <<FK>>
  + carryoverAmount  : DECIMAL(15,2)
  + reason           : TEXT
  + approvedById     : VARCHAR(10) <<FK>>
  + createdAt        : Timestamp
}

note right of MonthlyLoanAccrual
  Formule encours glissant :
  ──────────────────────────
  interestAccrued      = balanceAtMonthStart × 0.04
  balanceWithInterest  = balanceAtMonthStart × 1.04
  balanceAtMonthEnd    = balanceWithInterest
                         - repaymentReceived

  Décomposition remboursement R :
  ──────────────────────────────
  Si R >= interestAccrued :
    interestPart  = interestAccrued
    principalPart = R - interestAccrued
  Si R < interestAccrued :
    interestPart  = R
    principalPart = 0
    (capital inchangé, intérêts capitalisés)
end note

note bottom of LoanAccount
  Critères éligibilité :
  savings >= 100 000
  amount <= savings × 5
  amount ∈ [10 000, 1 000 000]
  activeLoans < 2
  dueBeforeDate < cassationDate

  Imputation multi-prêts :
  FIFO par requestedAt ASC
end note

LoanAccount "1" *-- "0..*" MonthlyLoanAccrual  : accrédite >
LoanAccount "1" *-- "0..*" LoanRepayment       : reçoit >
LoanAccount "1" ..> "0..1" CarryoverLoanRecord : génère à cassation >
LoanAccount --> LoanStatus

@enduml
\end{lstlisting}
\end{diagramBox}

% -----------------------------------------------------------
\section{D08 — Caisse de Secours}
% -----------------------------------------------------------

\begin{diagramBox}{D08 — Caisse de Secours}{Diagramme de Classes}
\begin{lstlisting}
@startuml D08_Secours
!theme plain
skinparam classAttributeIconSize 0
skinparam classFontSize 11
skinparam classBackgroundColor #EEF4FF
skinparam classBorderColor #004782
skinparam arrowColor #004782
title <b>Module 08 — Caisse de Secours</b>

class RescueFundLedger {
  + id               : VARCHAR(10) <<PK>>
  + fiscalYearId     : VARCHAR(10) <<FK, UNIQUE>>
  + totalBalance     : DECIMAL(15,2) = 0
  + targetPerMember  : DECIMAL(15,2) = 50000
  + minimumPerMember : DECIMAL(15,2) = 25000
  + memberCount      : INT
  + targetTotal      : DECIMAL(15,2)
  ' = targetPerMember × memberCount
}

class RescueFundEvent {
  + id             : VARCHAR(10) <<PK>>
  + ledgerId       : VARCHAR(10) <<FK>>
  + beneficiaryId  : VARCHAR(10) <<FK>>
  + eventType      : RescueEventType
  + amount         : DECIMAL(15,2)
  + authorizedById : VARCHAR(10) <<FK>>
  + eventDate      : DATE
  + description    : TEXT
  + createdAt      : Timestamp
}

class RescueFundPosition {
  + id            : VARCHAR(10) <<PK>>
  + membershipId  : VARCHAR(10) <<FK, UNIQUE>>
  + ledgerId      : VARCHAR(10) <<FK>>
  + paidAmount    : DECIMAL(15,2) = 0
  + balance       : DECIMAL(15,2) = 0
  + refillDebt    : DECIMAL(15,2) = 0
  + lastUpdatedAt : Timestamp
}

note right of RescueFundEvent
  RES-01 : Après paiement
    totalBalance >= minimumPerMember × memberCount
  RES-02 : amount = RescueEventAmount[eventType]
           (montant fixe, non négociable)
  RES-03 : authorizedById.role ∈
    {PRESIDENT, VICE_PRESIDENT}

  Calcul renflouement :
  déficit = targetTotal - totalBalance
  refillPerMember = déficit / memberCount
end note

RescueFundLedger "1" *-- "0..*" RescueFundEvent    : enregistre >
RescueFundLedger "1" *-- "0..*" RescueFundPosition : suit >

@enduml
\end{lstlisting}
\end{diagramBox}

% -----------------------------------------------------------
\section{D09 — Bénéficiaires}
% -----------------------------------------------------------

\begin{diagramBox}{D09 — Bénéficiaires}{Diagramme de Classes}
\begin{lstlisting}
@startuml D09_Beneficiaires
!theme plain
skinparam classAttributeIconSize 0
skinparam classFontSize 11
skinparam classBackgroundColor #EEF4FF
skinparam classBorderColor #004782
skinparam arrowColor #004782
title <b>Module 09 — Bénéficiaires Mensuels (Bouffe de la Tontine)</b>

enum BeneficiaryStatus {
  UNASSIGNED
  ASSIGNED
  DELIVERED
}

class BeneficiarySchedule {
  + id            : VARCHAR(10) <<PK>>
  + fiscalYearId  : VARCHAR(10) <<FK, UNIQUE>>
  + createdAt     : Timestamp
}

class BeneficiarySlot {
  + id              : VARCHAR(10) <<PK>>
  + scheduleId      : VARCHAR(10) <<FK>>
  + sessionId       : VARCHAR(10) <<FK>>
  + month           : TINYINT
  + slotIndex       : TINYINT
  + membershipId    : VARCHAR(10) <<FK>>
  + amountDelivered : DECIMAL(15,2) = 0
  + designatedById  : VARCHAR(10) <<FK>>
  + designatedAt    : Timestamp
  + status          : BeneficiaryStatus = UNASSIGNED
  + deliveredAt     : Timestamp
  + notes           : TEXT
  --
  + assign(memberId, byId) : void
  + deliver(amount)        : void
  <<UNIQUE(sessionId, slotIndex)>>
  <<UNIQUE(sessionId, membershipId)>>
}

note right of BeneficiarySlot
  Désignation manuelle :
  designatedById.role ∈
    {PRESIDENT, VICE_PRESIDENT, SUPER_ADMIN}

  amountDelivered : saisi manuellement
  par le Président — pas de calcul auto

  Un membre ne peut pas avoir
  2 slots dans la même session

  status = DELIVERED seulement si
  session.status = CLOSED

  slotIndex : permet plusieurs
  bénéficiaires par session
  (0, 1, 2 ...)
end note

BeneficiarySchedule "1" *-- "0..*" BeneficiarySlot : planifie >
BeneficiarySlot --> BeneficiaryStatus

@enduml
\end{lstlisting}
\end{diagramBox}

% -----------------------------------------------------------
\section{D10 — Cassation et Audit}
% -----------------------------------------------------------

\begin{diagramBox}{D10 — Cassation \& Audit}{Diagramme de Classes}
\begin{lstlisting}
@startuml D10_Cassation
!theme plain
skinparam classAttributeIconSize 0
skinparam classFontSize 10
skinparam classBackgroundColor #EEF4FF
skinparam classBorderColor #004782
skinparam arrowColor #004782
title <b>Module 10 — Cassation, Rollover & Audit</b>

class CassationRecord {
  + id                    : VARCHAR(10) <<PK>>
  + fiscalYearId          : VARCHAR(10) <<FK, UNIQUE>>
  + executedAt            : Timestamp
  + executedById          : VARCHAR(10) <<FK>>
  + totalSavingsReturned  : DECIMAL(15,2)
  + totalInterestReturned : DECIMAL(15,2)
  + totalDistributed      : DECIMAL(15,2)
  + memberCount           : INT
  + carryoverCount        : INT
  + notes                 : TEXT
}

class CassationRedistribution {
  + id             : VARCHAR(10) <<PK>>
  + cassationId    : VARCHAR(10) <<FK>>
  + membershipId   : VARCHAR(10) <<FK>>
  + savingsAmount  : DECIMAL(15,2)
  + interestAmount : DECIMAL(15,2)
  + totalReturned  : DECIMAL(15,2)
  --
  <<UNIQUE(cassationId, membershipId)>>
}

class PoolParticipantCassationShare {
  + id              : VARCHAR(10) <<PK>>
  + cassationId     : VARCHAR(10) <<FK>>
  + participantType : PoolParticipantType
  + principalAmount : DECIMAL(15,2)
  + interestEarned  : DECIMAL(15,2)
  + totalDistributed: DECIMAL(15,2)
  --
  <<UNIQUE(cassationId, participantType)>>
}

class AuditLog <<PartitionedByYear>> {
  + id          : VARCHAR(10) <<PK>>
  + actorId     : VARCHAR(10) <<FK>>
  + action      : VARCHAR(100)
  + entityType  : VARCHAR(50)
  + entityId    : VARCHAR(10)
  + beforeData  : JSON
  + afterData   : JSON
  + reason      : TEXT
  + ipAddress   : VARCHAR(45)
  + createdAt   : Timestamp
  + createdYear : YEAR <<VIRTUAL STORED>>
  --
  <<APPEND-ONLY>>
  <<PARTITION BY RANGE(createdYear)>>
}

note right of CassationRecord
  Séquence cassation (ordonnée) :
  ────────────────────────────────
  1. Identifier prêts actifs
     → CarryoverLoanRecord × prêt
  2. Redistribution intérêts finale
  3. CassationRedistribution par membre
  4. PoolParticipantCassationShare
  5. FiscalYear.status = CLOSED
  6. INSERT CassationRecord
  7. Notifier PRESIDENT + TRESORIER
  8. Créer FiscalYear N+1 (PENDING)
end note

CassationRecord "1" *-- "0..*" CassationRedistribution      : détaille >
CassationRecord "1" *-- "0..*" PoolParticipantCassationShare : distribue >

@enduml
\end{lstlisting}
\end{diagramBox}

% -----------------------------------------------------------
\section{D11 — Vue Globale des Agrégats}
% -----------------------------------------------------------

\begin{diagramBox}{D11 — Vue Globale}{Relations Inter-Modules}
\begin{lstlisting}
@startuml D11_Vue_Globale
!theme plain
skinparam classAttributeIconSize 0
skinparam classFontSize 9
skinparam classBackgroundColor #EEF4FF
skinparam classBorderColor #004782
skinparam arrowColor #004782
skinparam packageBackgroundColor #F4F8FF
skinparam packageBorderColor #7799CC
title <b>CAYA — Vue Globale des Relations Inter-Modules</b>

package "Identité" {
  class User
  class MemberProfile
}

package "Exercice" {
  class FiscalYear <<Root>>
  class FiscalYearConfig
  class Membership
  class ShareCommitment
}

package "Sessions" {
  class MonthlySession
  class SessionEntry
  class TransactionSequence
}

package "Épargne" {
  class SavingsLedger
  class InterestDistributionSnapshot
  class PoolParticipant
}

package "Prêts" {
  class LoanAccount <<Root>>
  class MonthlyLoanAccrual
  class CarryoverLoanRecord
}

package "Secours" {
  class RescueFundLedger
  class RescueFundEvent
  class RescueFundPosition
}

package "Bénéficiaires" {
  class BeneficiarySchedule
  class BeneficiarySlot
}

package "Cassation" {
  class CassationRecord
  class CassationRedistribution
  class PoolParticipantCassationShare
}

' Relations principales
User          "1" -- "1"    MemberProfile
MemberProfile "1" -- "0..*" Membership
FiscalYear    "1" *-- "0..*" Membership
FiscalYear    "1" -- "1"    FiscalYearConfig
FiscalYear    "1" *-- "0..*" MonthlySession
FiscalYear    "1" *-- "0..*" LoanAccount
FiscalYear    "1" -- "1"    RescueFundLedger
FiscalYear    "1" -- "1"    BeneficiarySchedule
FiscalYear    "1" -- "0..1" CassationRecord

Membership    "1" -- "1"    ShareCommitment
Membership    "1" -- "1"    SavingsLedger
Membership    "1" -- "1"    RescueFundPosition
Membership    "1" -- "0..*" LoanAccount

MonthlySession "1" *-- "0..*" SessionEntry
MonthlySession "1" -- "0..1" InterestDistributionSnapshot
MonthlySession "1" -- "0..*" BeneficiarySlot

LoanAccount   "1" -- "0..*" MonthlyLoanAccrual
LoanAccount   "1" -- "0..1" CarryoverLoanRecord

CassationRecord "1" -- "0..*" CassationRedistribution
CassationRecord "1" -- "0..*" PoolParticipantCassationShare

@enduml
\end{lstlisting}
\end{diagramBox}

% -----------------------------------------------------------
\section{D12 — Cas d'Utilisation}
% -----------------------------------------------------------

\begin{diagramBox}{D12 — Cas d'Utilisation}{Use Case Diagram}
\begin{lstlisting}
@startuml D12_UseCase
!theme plain
skinparam actorBackgroundColor #EEF4FF
skinparam actorBorderColor #004782
skinparam usecaseBackgroundColor #EEF4FF
skinparam usecaseBorderColor #004782
skinparam arrowColor #004782
left to right direction
title <b>CAYA — Diagramme des Cas d'Utilisation</b>

actor "Membre"           as MEMBRE
actor "Trésorier"        as TRES
actor "Président / VP"   as PRES
actor "Secrétaire"       as SEC
actor "Super Admin"      as SA

rectangle "Application CAYA" {

  package "Authentification" {
    usecase "Se connecter"          as UC_LOGIN
    usecase "Renouveler session"    as UC_REFRESH
  }

  package "Consultation Membre" {
    usecase "Voir sa fiche"         as UC_MY_ACCOUNT
    usecase "Voir historique"       as UC_HISTORY
    usecase "Voir ses prêts"        as UC_MY_LOANS
    usecase "Recevoir notifications" as UC_NOTIF
  }

  package "Saisie des Transactions" {
    usecase "Ouvrir session"        as UC_OPEN_SES
    usecase "Saisir transactions"   as UC_SAISIE
    usecase "Remboursement OOS"     as UC_OOS
    usecase "Clore session"         as UC_CLOSE_SES
    usecase "Distribuer intérêts"   as UC_DISTRIB
  }

  package "Gestion des Prêts" {
    usecase "Demander un prêt"      as UC_REQ_LOAN
    usecase "Autoriser un prêt"     as UC_AUTH_LOAN
    usecase "Débloquer un prêt"     as UC_DISB_LOAN
  }

  package "Gouvernance" {
    usecase "Désigner bénéficiaire" as UC_BENEF
    usecase "Autoriser secours"     as UC_SECOURS
    usecase "Configurer exercice"   as UC_CONFIG
    usecase "Consulter audit log"   as UC_AUDIT
    usecase "Exécuter cassation"    as UC_CASSATION
  }

  package "Rapports" {
    usecase "Exporter PDF"          as UC_PDF
    usecase "Exporter Excel"        as UC_XLS
  }
}

MEMBRE --> UC_LOGIN
MEMBRE --> UC_MY_ACCOUNT
MEMBRE --> UC_HISTORY
MEMBRE --> UC_MY_LOANS
MEMBRE --> UC_NOTIF

TRES   --> UC_LOGIN
TRES   --> UC_OPEN_SES
TRES   --> UC_SAISIE
TRES   --> UC_OOS
TRES   --> UC_CLOSE_SES
TRES   --> UC_DISTRIB
TRES   --> UC_DISB_LOAN
TRES   --> UC_CASSATION
TRES   --> UC_PDF
TRES   --> UC_XLS

PRES   --> UC_LOGIN
PRES   --> UC_AUTH_LOAN
PRES   --> UC_BENEF
PRES   --> UC_SECOURS
PRES   --> UC_AUDIT
PRES   --> UC_PDF

SEC    --> UC_LOGIN
SEC    --> UC_AUDIT
SEC    --> UC_PDF

SA     --> UC_CONFIG
SA     --> UC_AUDIT
SA     --> UC_CASSATION

@enduml
\end{lstlisting}
\end{diagramBox}

% -----------------------------------------------------------
\section{D13 — Diagramme de Composants}
% -----------------------------------------------------------

\begin{diagramBox}{D13 — Composants}{Component Diagram}
\begin{lstlisting}
@startuml D13_Composants
!theme plain
skinparam componentBackgroundColor #EEF4FF
skinparam componentBorderColor #004782
skinparam arrowColor #004782
skinparam interfaceBackgroundColor #FFF8E0
skinparam databaseBackgroundColor #E0F0FF
title <b>CAYA — Architecture des Composants</b>

package "Frontend Web (Next.js)" {
  [Dashboard]       as DASH
  [SessionSaisie]   as SAISIE
  [RapportsExport]  as RAPPORTS
  [ConfigPanel]     as CONFIG
}

package "Mobile (Flutter)" {
  [MemberApp]       as FLUTTER
}

package "Backend NestJS" {
  [AuthModule]      as AUTH
  [MembersModule]   as MEMB
  [SessionsModule]  as SESS
  [SavingsModule]   as SAV
  [LoansModule]     as LOAN
  [RescueModule]    as RES
  [CassationModule] as CASS
  [ReportsModule]   as REP
  [NotifModule]     as NOTIF
}

package "Infrastructure" {
  database "MySQL 8.0\ncaya_db" as DB
  [PrismaService]   as PRISMA
  [WhatsApp API]    as WA
  [SMS Provider]    as SMS
}

' Frontend → Backend
DASH     --> AUTH    : JWT Bearer
DASH     --> MEMB
DASH     --> SESS
SAISIE   --> SESS
RAPPORTS --> REP
CONFIG   --> AUTH

' Mobile → Backend
FLUTTER  --> AUTH    : JWT Bearer
FLUTTER  --> MEMB
FLUTTER  --> LOAN
FLUTTER  --> SAV

' Backend → Infrastructure
AUTH     --> PRISMA
MEMB     --> PRISMA
SESS     --> PRISMA
SAV      --> PRISMA
LOAN     --> PRISMA
RES      --> PRISMA
CASS     --> PRISMA
REP      --> PRISMA

PRISMA   --> DB

NOTIF    --> WA
NOTIF    --> SMS
CASS     --> NOTIF
SESS     --> NOTIF

@enduml
\end{lstlisting}
\end{diagramBox}

% -----------------------------------------------------------
\section{D14 — Diagramme de Déploiement}
% -----------------------------------------------------------

\begin{diagramBox}{D14 — Déploiement}{Deployment Diagram}
\begin{lstlisting}
@startuml D14_Deploiement
!theme plain
skinparam nodeBackgroundColor #EEF4FF
skinparam nodeBorderColor #004782
skinparam arrowColor #004782
title <b>CAYA — Diagramme de Déploiement</b>

node "Navigateur Web\n(Bureau / Trésorier)" as BROWSER {
  artifact "Next.js App\n(SPA / SSR)" as NEXT
}

node "Smartphone\n(Membres)" as PHONE {
  artifact "Flutter App\n(Android / iOS)" as FLUTTER
}

node "Serveur d'Application\n(VPS / Cloud)" as SERVER {
  artifact "NestJS API\n:3000" as API
  artifact "Prisma Client"     as PRISMA
}

node "Serveur Base de Données\n(MySQL 8.0)" as DBSERVER {
  database "caya_db\n(InnoDB)" as DB
}

node "Services Tiers" as EXTERNAL {
  artifact "WhatsApp\nBusiness API" as WA
  artifact "SMS Provider\n(Orange/MTN)" as SMS
}

BROWSER  --> SERVER   : HTTPS :443 / REST JSON
PHONE    --> SERVER   : HTTPS :443 / REST JSON
SERVER   --> DBSERVER : TCP :3306
API      --> PRISMA
PRISMA   --> DB
API      --> WA       : HTTPS Webhook
API      --> SMS       : HTTPS API

note right of SERVER
  - Node.js 20 LTS
  - NestJS 10
  - JWT RS256
  - Rate limiting : 100 req/min
  - Logs : Winston → fichiers
end note

note right of DBSERVER
  - MySQL 8.0
  - InnoDB Engine
  - utf8mb4_unicode_ci
  - Backup journalier
  - AuditLog partitionné
end note

@enduml
\end{lstlisting}
\end{diagramBox}

% ============================================================
\chapter{Diagrammes de Machine d'États}
% ============================================================

\begin{diagramBox}{D15 — État : FiscalYear}{State Machine}
\begin{lstlisting}
@startuml D15_State_FiscalYear
!theme plain
skinparam stateBackgroundColor #EEF4FF
skinparam stateBorderColor #004782
skinparam arrowColor #004782
title <b>Machine d'États — FiscalYear</b>

[*] --> PENDING : <<create>>\nPRESIDENT / SUPER_ADMIN

PENDING --> ACTIVE : activate()\n[snapshot FiscalYearConfig]\n[créer 12 MonthlySession DRAFT]\n[créer SavingsLedger par membre]\n[créer RescueFundLedger]\n[créer BeneficiarySchedule]

ACTIVE --> CASSATION : openCassation()\n[cassationDate atteinte]

CASSATION --> CLOSED : close()\n[CarryoverLoanRecord créés]\n[CassationRecord créé]\n[redistribution finale exécutée]\n[notifications envoyées]

CLOSED --> ARCHIVED : archive()\n[données figées]

ARCHIVED --> [*]

note right of PENDING
  Modifiable :
  - startDate / endDate
  - cassationDate
  - loanDueDate
  - interestPoolMethod
  — avant activation seulement —
end note

note right of ACTIVE
  Immuable :
  FiscalYearConfig.shareUnitAmount
  FiscalYearConfig.loanMonthlyRate
  — sauf SUPER_ADMIN force —
end note

note right of CASSATION
  Actions obligatoires avant CLOSED :
  1. Vérifier prêts → CarryoverLoanRecord
  2. Distribution intérêts finale
  3. CassationRedistribution par membre
end note

@enduml
\end{lstlisting}
\end{diagramBox}

\begin{diagramBox}{D16 — État : MonthlySession}{State Machine}
\begin{lstlisting}
@startuml D16_State_Session
!theme plain
skinparam stateBackgroundColor #EEF4FF
skinparam stateBorderColor #004782
skinparam arrowColor #004782
title <b>Machine d'États — MonthlySession</b>

[*] --> DRAFT : <<auto-créée>>\n lors de FiscalYear.activate()

DRAFT --> OPEN : open(TRESORIER)\n[session[n-1].status = CLOSED\n OU n = 1]

OPEN --> REVIEWING : requestReview(TRESORIER)\n[saisie terminée]

REVIEWING --> CLOSED : close(PRESIDENT)\n[computeTotals()]\n[distributeInterests()]

CLOSED --> OPEN : reopen(SUPER_ADMIN only)\n+ AuditLog obligatoire\n+ raison requise

note right of OPEN
  État actif pour la saisie :
  SessionEntry créables
  uniquement en statut OPEN

  Trigger MySQL bloque
  toute insertion si ≠ OPEN
end note

note right of CLOSED
  État terminal normal.
  Totaux snapshot figés.
  Intérêts redistribués.
  Uniquement SUPER_ADMIN
  peut rouvrir + AuditLog.
end note

@enduml
\end{lstlisting}
\end{diagramBox}

\begin{diagramBox}{D17 — État : LoanAccount}{State Machine}
\begin{lstlisting}
@startuml D17_State_Loan
!theme plain
skinparam stateBackgroundColor #EEF4FF
skinparam stateBorderColor #004782
skinparam arrowColor #004782
title <b>Machine d'États — LoanAccount</b>

[*] --> PENDING : requestLoan(membre)\n[LoanEligibilitySpec.validate()]

PENDING --> ACTIVE : disburse(TRESORIER)\n[PRESIDENT autorisé]\n[disbursedAt = today]\n[outstandingBalance = principalAmount]

PENDING --> [*] : reject(PRESIDENT)\n+ motif requis

ACTIVE --> PARTIALLY_REPAID : applyRepayment(R)\n[outstandingBalance > 0]

ACTIVE --> CLOSED : applyRepayment(R)\n[outstandingBalance = 0]

PARTIALLY_REPAID --> PARTIALLY_REPAID : applyRepayment(R)\n[still > 0]

PARTIALLY_REPAID --> CLOSED : applyRepayment(R)\n[outstandingBalance = 0]

ACTIVE --> [*] : <<cassation>>\nCarryoverLoanRecord créé\nnotification envoyée

PARTIALLY_REPAID --> [*] : <<cassation>>\nCarryoverLoanRecord créé\nnotification envoyée

note right of PENDING
  LoanEligibilitySpec vérifie :
  - savings >= 100 000
  - amount <= savings × 5
  - amount ∈ [10 000, 1 000 000]
  - activeLoans < 2
  - dueBeforeDate < cassationDate
end note

note right of ACTIVE
  Chaque mois :
  MonthlyLoanAccrual calculé
  encours[m] = encours[m-1] × 1.04
               - remboursement[m]
end note

note bottom of CLOSED
  outstandingBalance = 0.00
  totalRepaid = principalAmount
               + totalInterestAccrued
end note

@enduml
\end{lstlisting}
\end{diagramBox}

\begin{diagramBox}{D18 — État : MemberStatus}{State Machine}
\begin{lstlisting}
@startuml D18_State_Member
!theme plain
skinparam stateBackgroundColor #EEF4FF
skinparam stateBorderColor #004782
skinparam arrowColor #004782
title <b>Machine d'États — MemberStatus (dans Membership)</b>

[*] --> ACTIVE : enroll(fiscalYear)\n[TRESORIER / BUREAU]

ACTIVE --> INACTIVE : setInactive(PRESIDENT)\n[membre ne participe plus\n temporairement]

ACTIVE --> SUSPENDED : suspend(PRESIDENT)\n[non-respect des règles\n cotisations en retard]

ACTIVE --> EXCLUDED : exclude(PRESIDENT)\n[décision définitive]

INACTIVE --> ACTIVE : reactivate(PRESIDENT)\n[caisses à jour]

SUSPENDED --> ACTIVE : lift(PRESIDENT)\n[dette soldée\n cotisations à jour]

SUSPENDED --> EXCLUDED : exclude(PRESIDENT)\n[suspension persistante]

EXCLUDED --> [*]

note right of SUSPENDED
  Conséquences suspension :
  - Inéligible aux prêts
  - Inéligible au bénéfice
  - Épargne bloquée (pas de retrait)
  - Secours encore possible
end note

note right of EXCLUDED
  État terminal.
  Données conservées
  pour audit historique.
  Compte User désactivé.
end note

@enduml
\end{lstlisting}
\end{diagramBox}

% ============================================================
\chapter{Diagrammes de Séquence}
% ============================================================

\begin{diagramBox}{D19 — Séquence : Enregistrement Session Mensuelle}{Sequence Diagram}
\begin{lstlisting}
@startuml D19_Seq_Session
!theme plain
skinparam sequenceArrowThickness 1.5
skinparam participantBackgroundColor #EEF4FF
skinparam participantBorderColor #004782
skinparam actorBackgroundColor #FFF8E0
skinparam arrowColor #004782
skinparam noteBackgroundColor #FFFFF0
skinparam noteBorderColor #CCCC00
title <b>Séquence — Enregistrement d'une Réunion Mensuelle</b>

actor      "Trésorier"               as TRES
participant "SessionController"       as SC
participant "SessionService"          as SS
participant "TransactionRefFactory"   as TRF
participant "LoanRepaymentAllocator"  as LRA
participant "SavingsService"          as SAV
participant "InterestDistribService"  as IDS
participant "UnitOfWork (Prisma.$tx)" as UOW
database    "MySQL"                   as DB

TRES -> SC  : POST /sessions/{id}/entries\n[entries[]]
SC   -> SS  : recordEntries(sessionId, entries)
SS   -> DB  : SELECT status FROM monthly_sessions
DB   --> SS : 'OPEN'

note right of SS
  Guard : session.status MUST = 'OPEN'
  Sinon → 422 SESSION_NOT_OPEN
end note

loop pour chaque entry
  SS -> TRF : generateRef(type, sessionNumber, fiscalYear)
  TRF -> DB : SELECT FOR UPDATE\nFROM transaction_sequences
  DB --> TRF : lastSequence
  TRF -> DB : UPDATE lastSequence = lastSequence + 1
  TRF --> SS : "CAYA-2026-03-COT-0042"

  alt type = REMBOURSEMENT (amount saisi)
    SS -> LRA : allocateRepayment(membershipId, amount)
    LRA -> DB : SELECT loan_accounts\nWHERE status ≠ CLOSED\nORDER BY requested_at ASC
    DB --> LRA : loans[] (FIFO)
    LRA -> LRA : décomposer(amount, interestDue, capital)
    LRA --> SS : [{loanId, interestPart, principalPart}]
  end
end

SS -> UOW : $transaction(async (tx) => {
  UOW -> DB : INSERT session_entries[]
  UOW -> DB : INSERT loan_repayments[]
  UOW -> DB : UPDATE loan_accounts\n  outstandingBalance, totalRepaid
  UOW -> DB : INSERT monthly_loan_accruals
  UOW -> DB : INSERT savings_entries [DEPOSIT]
  UOW -> DB : UPDATE savings_ledgers.balance
  UOW -> DB : UPDATE rescue_fund_positions (si SECOURS)
  UOW -> DB : UPDATE monthly_sessions totaux
  UOW -> DB : COMMIT
})

SS --> SC  : {createdEntries[], sessionTotals{}}
SC --> TRES : 201 Created

TRES -> SC  : POST /sessions/{id}/distribute-interests
SC   -> IDS : distribute(sessionId)
IDS  -> DB  : SELECT interestPoolMethod\nFROM fiscal_year_configs
IDS  -> IDS : getStrategy(method)
IDS  -> DB  : computePool() selon stratégie
IDS  -> DB  : SELECT savings balances\n+ pool participants
IDS  -> UOW : INSERT snapshot\nINSERT allocations[]\nUPDATE savings_ledgers\nUPDATE pool_participants
IDS --> SC  : snapshot {pool, yield%, allocations[]}
SC --> TRES : 200 OK

@enduml
\end{lstlisting}
\end{diagramBox}

\begin{diagramBox}{D20 — Séquence : Distribution Mensuelle des Intérêts}{Sequence Diagram}
\begin{lstlisting}
@startuml D20_Seq_Interets
!theme plain
skinparam sequenceArrowThickness 1.5
skinparam participantBackgroundColor #EEF4FF
skinparam participantBorderColor #004782
skinparam actorBackgroundColor #FFF8E0
skinparam arrowColor #004782
title <b>Séquence — Distribution Mensuelle des Intérêts</b>

actor       "Trésorier"                  as TRES
participant "InterestDistribService"      as IDS
participant "TheoreticalStrategy"         as THEO
participant "ActualStrategy"              as ACT
participant "UnitOfWork"                  as UOW
database    "MySQL"                       as DB

TRES -> IDS : distributeInterests(sessionId)
IDS  -> DB  : SELECT interest_pool_method\nFROM fiscal_year_configs

alt method = THEORETICAL
  IDS -> THEO : computePool(sessionId)
  THEO -> DB  : SELECT SUM(interest_accrued)\nFROM monthly_loan_accruals\nWHERE session_id = ?
  DB --> THEO : totalInterestPool
  THEO --> IDS : totalInterestPool

else method = ACTUAL
  IDS -> ACT  : computePool(sessionId)
  ACT  -> DB  : SELECT SUM(amount)\nFROM session_entries\nWHERE type='RBT_INTEREST'\nAND session_id = ?
  DB --> ACT  : totalInterestPool
  ACT --> IDS : totalInterestPool
end

IDS -> DB : SELECT sl.membership_id,\n  sl.balance AS member_balance\nFROM savings_ledgers sl\nJOIN memberships m ON ...\nWHERE m.fiscal_year_id = ?

IDS -> DB : SELECT pp.type, pp.current_balance\nFROM pool_participants\nWHERE fiscal_year_id = ?

IDS -> IDS : totalBase = Σ member_balances\n            + Σ pool_balances

note right of IDS
  Pour chaque membre i :
  allocationAmount[i] =
    balance[i] / totalBase
    × totalInterestPool

  Vérification :
  Σ allocationAmount = totalInterestPool ± 1
end note

IDS -> UOW : $transaction(async (tx) => {
  UOW -> DB : INSERT interest_distribution_snapshots
  UOW -> DB : INSERT interest_allocations[]\n  (par membre)
  UOW -> DB : INSERT savings_entries [INTEREST_CREDIT]\n  (par membre)
  UOW -> DB : UPDATE savings_ledgers\n  balance += allocationAmount\n  totalInterestReceived += allocationAmount
  UOW -> DB : UPDATE pool_participants\n  currentBalance += allocationAmount\n  totalInterestReceived += allocationAmount
  UOW -> DB : COMMIT
})

IDS --> TRES : {snapshot{pool, yield%, allocations[]}}

@enduml
\end{lstlisting}
\end{diagramBox}

\begin{diagramBox}{D21 — Séquence : Processus de Cassation}{Sequence Diagram}
\begin{lstlisting}
@startuml D21_Seq_Cassation
!theme plain
skinparam sequenceArrowThickness 1.5
skinparam participantBackgroundColor #EEF4FF
skinparam participantBorderColor #004782
skinparam actorBackgroundColor #FFF8E0
skinparam arrowColor #004782
title <b>Séquence — Processus de Cassation (Clôture Exercice)</b>

actor       "Trésorier"               as TRES
participant "CassationService"         as CS
participant "InterestDistribService"   as IDS
participant "NotificationService"      as NOTIF
participant "UnitOfWork"               as UOW
database    "MySQL"                    as DB

TRES -> CS  : executeCassation(fiscalYearId)
CS   -> DB  : SELECT * FROM fiscal_years\nWHERE id = ? AND status = 'CASSATION'

note right of CS
  Étape 1 : Traiter prêts non soldés
end note

CS -> DB    : SELECT * FROM loan_accounts\nWHERE fiscal_year_id = ?\nAND status NOT IN ('CLOSED')

loop pour chaque prêt actif
  CS -> DB  : INSERT carryover_loan_records\n  carryoverAmount = outstandingBalance × 1.04
end

note right of CS
  Étape 2 : Distribution intérêts finale
end note

CS -> IDS   : distributeInterests(lastSessionId)
IDS --> CS  : snapshot final

note right of CS
  Étape 3 : Redistribution membres
end note

CS -> DB    : SELECT sl.membership_id,\n  sl.balance + sl.total_interest_received\nFROM savings_ledgers sl

CS -> UOW   : $transaction(async (tx) => {

  UOW -> DB : INSERT cassation_record\n  {totalDistributed, memberCount...}

  loop pour chaque membre
    UOW -> DB : INSERT cassation_redistributions\n  {savingsAmount, interestAmount, totalReturned}
  end

  loop RESCUE_FUND + BUREAU
    UOW -> DB : INSERT pool_participant_cassation_shares
  end

  UOW -> DB   : UPDATE fiscal_years\n  SET status = 'CLOSED'\n  closedAt = NOW()

  UOW -> DB   : COMMIT
})

note right of CS
  Étape 4 : Notifications + N+1
end note

CS -> NOTIF : notifyCassationComplete(\n  PRESIDENT, TRESORIER,\n  totalDistributed,\n  carryoverLoansCount\n)
NOTIF -> DB : LOG notification_sent

CS -> DB    : INSERT fiscal_years\n  {label:'2026-2027', status:'PENDING', ...}

CS --> TRES : {cassationRecord, redistributions[], carryovers[]}

@enduml
\end{lstlisting}
\end{diagramBox}

% ============================================================
\chapter*{Annexe — Commandes de Rendu}
\addcontentsline{toc}{chapter}{Annexe — Commandes de Rendu}
% ============================================================

\begin{noteBox}{Générer tous les diagrammes en ligne de commande}
\begin{lstlisting}[language=bash, numbers=none]
# Installer PlantUML (Java requis)
# https://plantuml.com/download

# Générer un diagramme unique
java -jar plantuml.jar D01_Auth_RBAC.puml

# Générer tous les diagrammes d'un dossier
java -jar plantuml.jar -o ./output/ *.puml

# Générer en SVG (meilleure qualité vectorielle)
java -jar plantuml.jar -tsvg *.puml

# Extension VSCode : PlantUML (jebbs.plantuml)
# Raccourci : Alt+D  (prévisualisation live)
# Raccourci : Ctrl+Shift+P > PlantUML: Export

# Site en ligne (copier-coller le code)
# https://www.plantuml.com/plantuml/uml/
\end{lstlisting}
\end{noteBox}

\begin{center}
\rowcolors{2}{cayaLight}{white}
\begin{tabular}{C{0.8cm}L{4.5cm}L{3cm}C{3cm}}
  \toprule
  \rowcolor{cayaBlue}
  \textcolor{white}{\textbf{\#}} &
  \textcolor{white}{\textbf{Diagramme}} &
  \textcolor{white}{\textbf{Type}} &
  \textcolor{white}{\textbf{Statut}} \\
  \midrule
  D01 & Auth \& RBAC              & Classes    & \cmark Complet \\
  D02 & Configuration             & Classes    & \cmark Complet \\
  D03 & Membres                   & Classes    & \cmark Complet \\
  D04 & Exercice Fiscal           & Classes    & \cmark Complet \\
  D05 & Sessions \& Transactions  & Classes    & \cmark Complet \\
  D06 & Épargne \& Intérêts       & Classes    & \cmark Complet \\
  D07 & Prêts                     & Classes    & \cmark Complet \\
  D08 & Caisse de Secours         & Classes    & \cmark Complet \\
  D09 & Bénéficiaires             & Classes    & \cmark Complet \\
  D10 & Cassation \& Audit        & Classes    & \cmark Complet \\
  D11 & Vue Globale               & Classes    & \cmark Complet \\
  D12 & Cas d'Utilisation         & Use Case   & \cmark Complet \\
  D13 & Composants                & Composants & \cmark Complet \\
  D14 & Déploiement               & Déploiement& \cmark Complet \\
  D15 & État — FiscalYear         & États      & \cmark Complet \\
  D16 & État — MonthlySession     & États      & \cmark Complet \\
  D17 & État — LoanAccount        & États      & \cmark Complet \\
  D18 & État — MemberStatus       & États      & \cmark Complet \\
  D19 & Séquence — Session        & Séquence   & \cmark Complet \\
  D20 & Séquence — Intérêts       & Séquence   & \cmark Complet \\
  D21 & Séquence — Cassation      & Séquence   & \cmark Complet \\
  \bottomrule
\end{tabular}
\end{center}

\vfill
\begin{center}
  \textcolor{cayaBlue}{\rule{0.5\textwidth}{0.4pt}}\\[0.3cm]
  {\small\textit{Cahier des Diagrammes UML — CAYA}}\\
  {\small\textcolor{gray}{21 diagrammes · PlantUML · Version 1.0 · Mars 2026}}
\end{center}

\end{document}
